% Source: https://artofproblemsolving.com/community/c3h2915810_2014_emcc_guts_round_p13_p24__exeter_math_club_competition
% Stable identity: aops-topic-2915810-q17
The answer is $\boxed5$. Here is an exhaustive recursive counting proof: the partner of vertex 0 must have odd index; its chord splits the remaining vertices into two even consecutive lists, whose noncrossing matchings are chosen independently. This generates each noncrossing matching exactly once (42 total). For every triple of chords, sort its six endpoints around the circle and test all cyclic starting positions; the prohibited pattern occurs exactly when the three chords pair consecutive endpoint pairs in some such order. The following complete finite computation leaves five matchings.
\begin{verbatim}
from functools import lru_cache
from itertools import combinations
@lru_cache(None)
def match(v):
 if not v:return ((),)
 out=[]
 for j in range(1,len(v),2):
  for left in match(v[1:j]):
   for right in match(v[j+1:]):
    out.append(((v[0],v[j]),)+left+right)
 return tuple(out)
def forbidden(m):
 for es in combinations(m,3):
  pts=sorted(x for e in es for x in e)
  target={frozenset(e) for e in es}
  for k in range(6):
   q=pts[k:]+pts[:k]
   if {frozenset((q[i],q[i+1])) for i in (0,2,4)}==target:
    return True
 return False
all_m=match(tuple(range(10)))
assert len(all_m)==42
assert sum(not forbidden(m) for m in all_m)==5
\end{verbatim}
\begin{asy}
usepackage("amsmath");
size(7cm);
pair[] V;for(int k=0;k<10;++k){V.push(3*dir(36*k));dot(V[k]);label(string(k),1.13*V[k]);}draw(circle((0,0),3),gray);for(int k=0;k<5;++k)draw(V[k]--V[9-k],blue);
\end{asy}
